\chapter{Gap Analysis}

This chapter systematically analyzes limitations in existing dropout prediction systems and justifies the proposed AHFS-TA framework.

\section{Limitations of Existing Approaches}

\subsection{Gap 1: Narrow Feature Engineering}

\textbf{Observation}: Most studies (92\% of reviewed literature) rely exclusively on structured administrative data: demographics, grades, financial records \cite{delen2010comparative, huang2020deep, adnan2021student}.

\textbf{Limitation}: Psychosocial factors (student sentiment, engagement quality, cognitive struggles) known to influence dropout \cite{tinto2012completing} remain unquantified.

\textbf{Consequence}: Models miss critical behavioral and emotional indicators extractable from textual interactions (forum posts, feedback, emails).

\textbf{Solution Needed}: Multimodal learning integrating structured data with LLM-derived psychosocial features.

\subsection{Gap 2: Static Modeling Paradigm}

\textbf{Observation}: 85\% of existing work treats semesters as independent snapshots, ignoring temporal progression \cite{kotsiantis2003machine, delen2010comparative}.

\textbf{Limitation}: Performance trends (improving/declining GPA trajectories, critical transition periods) are stronger dropout indicators than single-semester metrics \cite{tinto1975dropout}.

\textbf{Consequence}: Models fail to capture semester-wise dynamics, momentum shifts, and cumulative effects.

\textbf{Solution Needed}: Temporal attention mechanisms modeling sequential patterns across semesters.

\subsection{Gap 3: Fixed Feature Sets}

\textbf{Observation}: All reviewed studies use manually selected or domain-expert-curated feature sets fixed before training.

\textbf{Limitation}: No adaptive selection during training to optimize feature subsets dynamically based on model feedback.

\textbf{Consequence}: Suboptimal performance, reduced interpretability, unnecessary computational cost from redundant features.

\textbf{Solution Needed}: Adaptive hierarchical selection combining multiple importance perspectives (SHAP, LLM attention, temporal significance).

\subsection{Gap 4: Black-Box Predictions}

\textbf{Observation}: Only 18\% of studies integrate explainability methods; most provide predictions without reasoning.

\textbf{Limitation}: Educators and administrators cannot trust or act on opaque predictions.

\textbf{Consequence}: Limited real-world deployment despite high accuracy claims; institutional resistance to black-box systems.

\textbf{Solution Needed}: Integrated explainable AI (SHAP, attention weights) generating transparent, actionable insights.

\subsection{Gap 5: Incomplete Benchmarking}

\textbf{Observation}: Most papers evaluate 1-3 models; comprehensive cross-paradigm comparisons (classical ML + ensemble + deep learning) rare.

\textbf{Limitation}: Unclear which modeling approach best suits educational data characteristics.

\textbf{Consequence}: Practitioners lack guidance on optimal algorithm selection for their institutional context.

\textbf{Solution Needed}: Rigorous evaluation of diverse models (single classifiers, ensembles, deep learning) with 10-fold cross-validation on identical data splits.

\section{Proposed Solution: AHFS-TA Framework}

To address identified gaps, we propose \textbf{Adaptive Hierarchical Feature Selection with Temporal Attention (AHFS-TA)}:

\subsection{Component 1: LLM-Based Psychosocial Feature Extraction}

\begin{itemize}
\item Uses DistilBERT to extract 4 features from student texts: Sentiment, Engagement, TopicConsistency, CognitiveLoad
\item Addresses Gap 1 (multimodal learning)
\item Expected contribution: +1.5-2\% accuracy improvement
\end{itemize}

\subsection{Component 2: Temporal Attention Network}

\begin{itemize}
\item Bidirectional GRU processes semester sequences
\item Multi-head attention identifies critical periods
\item Addresses Gap 2 (temporal modeling)
\item Expected contribution: +1-1.5\% accuracy improvement
\end{itemize}

\subsection{Component 3: Adaptive Hierarchical Feature Selection}

\begin{itemize}
\item Three-stream meta-ranking: SHAP + LLM attention + temporal significance
\item Dynamic selection at epoch 5 during training
\item Addresses Gap 3 (adaptive selection)
\item Expected contribution: +0.5-1\% accuracy + 20-30\% feature reduction
\end{itemize}

\subsection{Component 4: Integrated Explainability}

\begin{itemize}
\item SHAP values for feature attribution
\item Attention weights for temporal focus visualization
\item Addresses Gap 4 (black-box predictions)
\item Expected outcome: Interpretable predictions for educators
\end{itemize}

\subsection{Component 5: Comprehensive Benchmarking}

\begin{itemize}
\item Evaluate 6 baseline models: Decision Tree, Naive Bayes, Random Forest, AdaBoost, XGBoost, Neural Network
\item Rigorous 10-fold cross-validation
\item Addresses Gap 5 (incomplete evaluation)
\item Expected outcome: Identify best-performing paradigm for educational data
\end{itemize}

\section{Research Hypothesis}

\textbf{H1}: Multimodal learning (structured + LLM features) outperforms structured-only baselines by $\geq$ 1.5\%.

\textbf{H2}: Temporal attention modeling outperforms static models by $\geq$ 1\%.

\textbf{H3}: Adaptive feature selection reduces dimensionality by $\geq$ 20\% while maintaining or improving accuracy.

\textbf{H4}: AHFS-TA achieves $\geq$ 90\% accuracy and $\geq$ 92\% AUC-ROC, exceeding current SOTA (87.3\%, Liang 2022).

\section{Summary}

Existing dropout prediction systems exhibit five critical gaps: narrow feature engineering, static modeling, fixed feature sets, black-box predictions, and incomplete benchmarking. The proposed AHFS-TA framework addresses all gaps through multimodal LLM integration, temporal attention, adaptive selection, integrated explainability, and comprehensive evaluation. Next chapter details the methodology.
